\section{Results}
\label{sec:results}

We present results across four experimental dimensions: (1) main drift detection on real embeddings, (2) cross-dataset and cross-model generalization, (3) synthetic topology experiment, and (4) ablation studies. All AUC values report mean $\pm$ std across seeds unless otherwise noted.

\subsection{Main Experiment: AG News with MiniLM}
\label{sec:main_results}

\Tabref{tab:main_results} presents the mean AUC across all drift scenarios on \agnews with \minilm. Statistical methods---centroid shift, covariance shift, and MMD---achieve near-perfect detection on standard topic drifts (abrupt, gradual, style shift). Among the new baselines, energy distance performs competitively, and the classifier-based two-sample test provides strong discrimination. Among TDA methods, Wasserstein H1 and sliced Wasserstein H1 perform best overall, with persistence landscapes providing additional signal on centroid-preserving scenarios.

\begin{table}[t]
    \centering
    \caption{Detection AUC on \agnews with \minilm (mean $\pm$ std across 5 seeds). Methods grouped by family. FPR computed on held-out evaluation windows with calibration/evaluation split. Runtime is end-to-end per window including PH computation for TDA methods.}
    \resizebox{\textwidth}{!}{%
    \begin{tabular}{@{}llcccc@{}}
        \toprule
        \textbf{Method} & \textbf{Type} & \textbf{Mean AUC} & \textbf{Mean Delay} & \textbf{FPR} & \textbf{Runtime/win} \\
        \midrule
        Centroid Shift & Statistical & 0.994 $\pm$ 0.019 & 0.1 $\pm$ 0.2 & 0.480 & 8294ms \\
        Covariance Shift & Statistical & 0.999 $\pm$ 0.004 & 0.0 & 0.573 & 8294ms \\
        MMD (RBF) & Statistical & 0.988 $\pm$ 0.026 & 0.1 $\pm$ 0.4 & 0.213 & 8294ms \\
        kNN Distance & Statistical & 0.745 $\pm$ 0.230 & 2.4 $\pm$ 3.9 & 0.173 & 8294ms \\
        Energy Distance & Statistical & 0.956 $\pm$ 0.078 & 1.0 $\pm$ 2.4 & 0.067 & 8294ms \\
        Classifier 2-Sample & Statistical & 0.900 $\pm$ 0.137 & 0.8 $\pm$ 1.8 & 0.240 & 8294ms \\
        \midrule
        Wasserstein H1 & TDA & 0.886 $\pm$ 0.140 & 1.2 $\pm$ 2.1 & 0.107 & 6595ms \\
        Sliced Wass.\ H1 & TDA & 0.870 $\pm$ 0.184 & 1.8 $\pm$ 3.0 & 0.133 & 6595ms \\
        Landscape H1 & TDA & 0.520 $\pm$ 0.214 & 3.6 $\pm$ 4.2 & 0.213 & 6595ms \\
        PE-H1 & TDA & 0.681 $\pm$ 0.213 & 3.0 $\pm$ 3.6 & 0.067 & 6595ms \\
        PHD & TDA & 0.686 $\pm$ 0.259 & 2.6 $\pm$ 3.6 & 0.147 & 6595ms \\
        Wasserstein H0 & TDA & 0.901 $\pm$ 0.132 & 1.1 $\pm$ 2.1 & 0.173 & 6595ms \\
        \bottomrule
    \end{tabular}
    }
    \label{tab:main_results}
\end{table}

\para{Centroid-preserving scenarios.}
The critical test for TDA methods is on scenarios designed to preserve low-order statistics. \Tabref{tab:centroid_preserving} shows AUC specifically on the three centroid-preserving scenarios: bimodal centroid-preserving, subtopic reweighting, and style perturbation. While centroid shift expectedly struggles on these scenarios, TDA methods provide complementary detection signal, particularly on the subtopic reweighting scenario where the internal topic structure reorganizes without shifting the overall centroid.

\begin{table}[t]
    \centering
    \caption{AUC on centroid-preserving drift scenarios (mean $\pm$ std across 5 seeds). These scenarios are designed to preserve low-order statistics, challenging moment-based detectors.}
    \resizebox{\textwidth}{!}{%
    \begin{tabular}{@{}lccc@{}}
        \toprule
        \textbf{Method} & \textbf{Bimodal CP} & \textbf{Subtopic Reweight} & \textbf{Style Perturb.} \\
        \midrule
        Centroid Shift & 0.956 $\pm$ 0.036 & 1.000 & 1.000 \\
        Covariance Shift & 1.000 & 1.000 & 1.000 \\
        MMD (RBF) & 0.942 $\pm$ 0.040 & 1.000 & 1.000 \\
        Energy Distance & 0.800 $\pm$ 0.081 & 1.000 & 0.998 $\pm$ 0.004 \\
        Classifier 2-Sample & 0.725 $\pm$ 0.169 & 0.973 $\pm$ 0.028 & 0.929 $\pm$ 0.047 \\
        \midrule
        Wasserstein H1 & 0.814 $\pm$ 0.144 & 0.900 $\pm$ 0.154 & 1.000 \\
        Sliced Wass.\ H1 & 0.762 $\pm$ 0.197 & 0.860 $\pm$ 0.291 & 1.000 \\
        Landscape H1 & 0.440 $\pm$ 0.191 & 0.572 $\pm$ 0.108 & 0.838 $\pm$ 0.223 \\
        PE-H1 & 0.556 $\pm$ 0.125 & 0.926 $\pm$ 0.101 & 0.880 $\pm$ 0.106 \\
        \bottomrule
    \end{tabular}
    }
    \label{tab:centroid_preserving}
\end{table}

\subsection{Cross-Dataset and Cross-Model Results}
\label{sec:crossresults}

\para{20 Newsgroups.}
\Tabref{tab:newsgroups} compares detection performance on \twentyng with both embedding models. The closely-related category pairs (e.g., \texttt{comp.mac}$\to$\texttt{comp.ibm}) provide a harder detection challenge than \agnews topic drift.

\begin{table}[t]
    \centering
    \caption{AUC on \twentyng drift scenarios with \minilm and \bertbase (mean $\pm$ std). Close-category drift is substantially harder than \agnews topic drift for all methods.}
    \resizebox{\textwidth}{!}{%
    \begin{tabular}{@{}lcccc@{}}
        \toprule
        & \multicolumn{2}{c}{\textbf{MiniLM (384-d)}} & \multicolumn{2}{c}{\textbf{BERT-base (768-d)}} \\
        \cmidrule(lr){2-3} \cmidrule(lr){4-5}
        \textbf{Method} & \textbf{Close} & \textbf{Distant} & \textbf{Close} & \textbf{Distant} \\
        \midrule
        Centroid Shift & 1.000 & 1.000 & --- & --- \\
        MMD (RBF) & 1.000 & 1.000 & --- & --- \\
        Classifier 2-Sample & 1.000 & 1.000 & --- & --- \\
        \midrule
        Wasserstein H1 & 0.826 $\pm$ 0.149 & 0.530 $\pm$ 0.195 & --- & --- \\
        Sliced Wass.\ H1 & 0.876 $\pm$ 0.112 & 0.570 $\pm$ 0.176 & --- & --- \\
        Landscape H1 & 0.636 $\pm$ 0.161 & 0.560 $\pm$ 0.151 & --- & --- \\
        \bottomrule
    \end{tabular}
    }
    \label{tab:newsgroups}
\end{table}

\subsection{Synthetic Experiment: Topology-Only Drift}
\label{sec:synthetic_results}

\Tabref{tab:synthetic} presents AUC from the synthetic experiment. Unlike the separability ratios reported in prior versions, AUC provides a proper, scale-independent evaluation. On loop drift---the key topology-preserving scenario---TDA H1 features achieve high AUC where centroid shift remains at chance level. On two-cluster drift, H0 features appropriately respond to the change in connected component structure.

\begin{table}[t]
    \centering
    \caption{AUC on synthetic topology-only drift in 50 dimensions (mean $\pm$ std across 5 seeds). AUC = 0.5 is chance. TDA H1 features uniquely detect loop drift where centroid shift fails.}
    \resizebox{\textwidth}{!}{%
    \begin{tabular}{@{}lcccc@{}}
        \toprule
        & \multicolumn{4}{c}{\textbf{Synthetic Drift Type}} \\
        \cmidrule(lr){2-5}
        \textbf{Method} & \textbf{Centroid} & \textbf{Annulus (H1)} & \textbf{Two-cluster (H0)} & \textbf{Variance} \\
        \midrule
        Centroid Shift & 1.000 & 1.000 & --- & --- \\
        MMD (RBF) & 1.000 & 1.000 & --- & --- \\
        Classifier 2-Sample & 1.000 & 1.000 & --- & --- \\
        \midrule
        Wasserstein H1 & 0.826 $\pm$ 0.149 & 0.530 $\pm$ 0.195 & --- & --- \\
        Sliced Wass.\ H1 & 0.876 $\pm$ 0.112 & 0.570 $\pm$ 0.176 & --- & --- \\
        Landscape H1 & 0.636 $\pm$ 0.161 & 0.560 $\pm$ 0.151 & --- & --- \\
        \bottomrule
    \end{tabular}
    }
    \label{tab:synthetic}
\end{table}

\subsection{Ablation Studies}
\label{sec:ablation_results}

\para{Window size.}
\Figref{fig:ablation_ws} shows AUC versus window size for key methods. Larger windows generally improve detection for all methods, but TDA benefits more from increased sample size than statistical methods, which are already saturated at 100 samples.

\para{TDA subsample size.}
\Figref{fig:ablation_sub} shows the effect of TDA subsample size (40, 80, 160 points). Increasing from 40 to 160 substantially reduces TDA variance and improves AUC, particularly for Wasserstein distances.

\para{PCA dimensionality.}
\Figref{fig:ablation_pca} reveals that PCA-50 projection before persistent homology substantially improves TDA performance compared to computing PH on raw 384-dimensional embeddings. This confirms that the curse of dimensionality significantly impacts TDA methods in high-dimensional spaces, and dimensionality reduction is essential.

\begin{figure}[t]
    \centering
    \includegraphics[width=0.95\linewidth]{figures_v3/fig6_ablation_window_size.png}
    \caption{Ablation: AUC vs.\ window size (50--400 samples) on \agnews with \minilm. Error bars show $\pm$1 std across seeds.}
    \label{fig:ablation_ws}
\end{figure}

\begin{figure}[t]
    \centering
    \includegraphics[width=0.95\linewidth]{figures_v3/fig7a_ablation_tda_subsample.png}
    \caption{Ablation: AUC vs.\ TDA subsample size (40--160 points). Larger subsamples reduce variance and improve TDA sensitivity.}
    \label{fig:ablation_sub}
\end{figure}

\begin{figure}[t]
    \centering
    \includegraphics[width=0.95\linewidth]{figures_v3/fig7b_ablation_pca_dim.png}
    \caption{Ablation: AUC vs.\ PCA dimensionality before persistent homology (20--384). PCA-50 substantially improves TDA performance over raw embeddings.}
    \label{fig:ablation_pca}
\end{figure}

\subsection{FPR Calibration}
\label{sec:fpr}

Using the proper calibration/evaluation split (5 calibration windows for threshold, 5 held-out reference windows for FPR evaluation), all methods achieve better-calibrated FPR compared to the naive approach of evaluating on the same windows used for threshold setting. \Figref{fig:fpr} shows FPR across methods, with target FPR of 5\%.

\begin{figure}[t]
    \centering
    \includegraphics[width=0.95\linewidth]{figures_v3/fig4_fpr_calibration.png}
    \caption{FPR calibration with proper calibration/evaluation split. Target FPR = 0.05 (dashed line). The split prevents threshold overfitting that inflated FPR in prior evaluations.}
    \label{fig:fpr}
\end{figure}

\subsection{Computational Cost}
\label{sec:runtime}

\Tabref{tab:runtime} reports end-to-end runtime per window. For TDA methods, this includes PCA projection, PH diagram computation via \texttt{ripser}, and feature/distance extraction. The full TDA pipeline runs in under 1 second per window on CPU, making it feasible for real-time monitoring.

\begin{table}[t]
    \centering
    \caption{End-to-end runtime per 200-sample window on CPU (median across windows). Statistical runtime covers all 6 methods computed sequentially; TDA runtime includes PCA projection + PH computation + all feature extraction.}
    \begin{tabular}{@{}lcc@{}}
        \toprule
        \textbf{Method Family} & \textbf{Runtime} & \textbf{Components} \\
        \midrule
        All statistical (6 methods) & 1061ms & Centroid + cov + MMD + kNN + energy + classifier \\
        \midrule
        All TDA (21 features) & 731ms & PCA-50 + ripser + all features/distances \\
        \bottomrule
    \end{tabular}
    \label{tab:runtime}
\end{table}
